\section*{Hoofdstuk \ref{ch:b3:deformation-mechanisms} --- Vervormingsmechanismen en schuifzones}

\begin{solution}{exo:b3:deformation-mechanisms:1}
Forsteriet: $1573/2163 = 0.73$; ijs: $263/273 = 0.96$. Ijs ligt veel dichter
bij zijn smeltpunt en kruipt makkelijker: gletsjers stromen aan het oppervlak
van de Aarde.
\end{solution}

\begin{solution}{exo:b3:deformation-mechanisms:2}
$\gamma = \tan\ang{70} = 2.75$ en $\tan\ang{80} = 5.67$. Dicht bij \ang{90} verandert de
tangens heel snel met de hoek: een fout van één graad bij $\psi =
\ang{84}$ verandert $\gamma$ met een vijfde.
\end{solution}

\begin{solution}{exo:b3:deformation-mechanisms:3}
$\sigma = (10^{3.56}/10)^{1/1.26} = 363.1^{0.794} = \qty{108}{MPa}$.
\end{solution}

\begin{solution}{exo:b3:deformation-mechanisms:4}
$E/RT = 223000/(8.314\times773.15) = 34.69$, dus $\dot\varepsilon =
\num{1.1e-4}\times20^4\times e^{-34.69} = \num{1.1e-4}\times\num{1.6e5}\times
\num{8.6e-16} = \qty{1.5e-14}{s^{-1}}$.
\end{solution}

\begin{solution}{exo:b3:deformation-mechanisms:5}
$\eta = \sigma/(2\dot\varepsilon)$: $\num{65.4e6}/\num{2e-14} =
\qty{3.3e21}{Pa.s}$ bij \qty{400}{\celsius} en $\num{18.0e6}/\num{2e-14} =
\qty{9.0e20}{Pa.s}$ bij \qty{500}{\celsius}: de middelste korst bij deze
temperaturen is ongeveer even viskeus als de mantel onder de ijskappen.
\end{solution}

\begin{solution}{exo:b3:deformation-mechanisms:6}
$n = \log(\num{1.6e-4}/\num{e-5})/\log2 = \log16/\log2 = 4$: \omterm{def:b3:deformation-mechanisms:creep}{dislocatiekruip},
waarvan de \omterm{def:b3:deformation-mechanisms:flow-law}{spanningsexponent} ruim boven 1 ligt.
\end{solution}

\begin{solution}{exo:b3:deformation-mechanisms:7}
$\ln10 = (E/R)(1/1000 - 1/1100)$, dus $E = 8.314\times2.3026/\num{9.091e-5} =
\qty{211}{kJ.mol^{-1}}$.
\end{solution}

\begin{solution}{exo:b3:deformation-mechanisms:8}
Nabarro--Herring ($d^{-2}$): 100 keer sneller; Coble ($d^{-3}$): 1000 keer.
Langs de korrelgrenzen gebruiken de diffunderende atomen een weg waarvan de dwarsdoorsnede, $\delta
d$, krimpt met de korrelgrootte mee, bovenop de kortere afstand en het kleinere
volume dat verplaatst moet worden.
\end{solution}

\begin{solution}{exo:b3:deformation-mechanisms:9}
$\sigma_b^{2.5} = \num{3.59e-6}\,d^{-3}/\num{6.54e-17}$: voor $d = 1000$,
$\sigma_b = 54.9^{0.4} = \qty{5.0}{MPa}$; voor $d = 100$, $\sigma_b =
\num{54900}^{0.4} = \qty{79}{MPa}$. De helling is $-3/2.5 = -1.2$: een tien keer
fijnere korrel verhoogt de grenssterkte met $10^{1.2} = 16$.
\end{solution}

\begin{solution}{exo:b3:deformation-mechanisms:10}
Op beide paden is het drijvende concentratieverschil hetzelfde en de
afstand is $d$; het volume dat per eenheid rek verplaatst moet worden is $d^3$. Het rooster
biedt een dwarsdoorsnede $\propto d^2$, dus de snelheid gaat als $d^2/(d\,d^3) =
d^{-2}$; de korrelgrenzen bieden $\propto\delta d$, dus $\delta d/(d\,d^3) = \delta
d^{-3}$. De verhouding van de grens- en roostersnelheden is $(B_B\delta D_B)/(B_LD_Ld)$,
groter dan 1 wanneer $d < (B_B/B_L)\,\delta D_B/D_L$.
\end{solution}

\begin{solution}{exo:b3:deformation-mechanisms:11}
$\eta \propto \dot\varepsilon^{(1-n)/n} = \dot\varepsilon^{-2.5/3.5}$: tien keer
sneller deelt de viscositeit door $10^{0.714} = 5.2$. Een band die toevallig
sneller vervormt, wordt zwakker en neemt zo meer van de vervorming voor zijn rekening: de rek
verzamelt zich erin.
\end{solution}

\begin{solution}{exo:b3:deformation-mechanisms:12}
Vóór: $\dot\varepsilon_{\mathrm{dis}} = \num{6.54e-17}\times10^{3.5} =
\qty{2.07e-13}{s^{-1}}$ en $\dot\varepsilon_{\mathrm{dif}} = \num{3.59e-6}
\times10/3000^3 = \qty{1.3e-15}{s^{-1}}$: \qty{2.08e-13}{s^{-1}}, \omterm{def:b3:deformation-mechanisms:creep}{dislocatiekruip}.
Na: $\dot\varepsilon_{\mathrm{dif}} = \num{3.59e-6}\times10/30^3 =
\qty{1.33e-9}{s^{-1}}$, \omterm{def:b3:deformation-mechanisms:creep}{diffusiekruip}, \num{6400} keer sneller bij dezelfde
spanning: herkristallisatie heeft de zone enorm verzacht.
\end{solution}

\begin{solution}{pb:b3:deformation-mechanisms:1}
\textbf{1.} De gang draait van \ang{90} naar \ang{6} ten opzichte van de \omterm{def:b3:deformation-mechanisms:shear-zone}{schuifzone}: $\psi =
\ang{84}$ en $\gamma = \tan\ang{84} = 9.51$.

\textbf{2.} $\gamma w = 9.51\times400 = \qty{3804}{m}$, ongeveer \qty{3.8}{km}.

\textbf{3.} $S_1 = \sqrt{1 + 9.51^2/4} + 9.51/2 = 4.859 + 4.755 = 9.614$,
$S_2 = 1/S_1 = 0.104$ en $S_1/S_2 = 92.4$.

\textbf{4.} $\tan2\theta' = 2/9.51 = 0.210$, $\theta' = \ang{5.9}$: de
foliatie, dicht bij de lange as, draait steeds dichter naar de \omterm{def:b3:deformation-mechanisms:shear-zone}{schuifzone} naarmate $\gamma$
groeit.

\textbf{5.} Noordwaarts gevolgd wordt de gang in de \omterm{def:b3:deformation-mechanisms:shear-zone}{schuifzone} naar het oosten meegevoerd en
verlaat de \omterm{def:b3:deformation-mechanisms:shear-zone}{schuifzone} verder oostelijk: de noordelijke wand is ten opzichte van de zuidelijke naar het oosten opgeschoven.
Vanaf de zuidelijke wand gezien is de verre kant naar rechts verschoven: de \omterm{def:b3:deformation-mechanisms:shear-zone}{schuifzone} is
dextraal.

\textbf{6.} In het vlak dat de rekkingslineatie bevat, hier horizontaal voor een zone met
horizontale verschuiving, en loodrecht op de verticale foliatie: een
horizontale doorsnede. De S-foliatie helt tegen de klok in van de \omterm{def:b3:deformation-mechanisms:shear-zone}{schuifzone} af, richting
ONO, zoals de lange as van de rekellips.

\textbf{7.} $\sigma = (10^{3.56}/20)^{1/1.26} = 181.5^{0.794} = \qty{62}{MPa}$.

\textbf{8.} $242.1^{0.794} = \qty{78}{MPa}$ voor \qty{15}{\micro m} en
$145.2^{0.794} = \qty{52}{MPa}$ voor \qty{25}{\micro m}.

\textbf{9.} Ja: de experimenten besloegen van 34 tot \qty{268}{MPa}.

\textbf{10.} De zone vervormde bij \qty{450}{\celsius}, ver onder de 800 tot
\qty{1100}{\celsius} van de experimenten; een verband dat niet van de
temperatuur afhangt, kan tot daar worden doorgetrokken, al blijft het
een extrapolatie.

\textbf{11.} $D = 3631/100^{1.26} = 3631/331 = \qty{11}{\micro m}$.

\textbf{12.} $E/RT = 223000/(8.314\times723.15) = 37.09$ en $e^{-37.09} =
\num{7.79e-17}$.

\textbf{13.} $\dot\varepsilon = \num{1.1e-4}\times62^4\times\num{7.79e-17} =
\num{1.1e-4}\times\num{1.478e7}\times\num{7.79e-17} = \qty{1.27e-13}{s^{-1}}$.

\textbf{14.} $\dot\gamma = 2\dot\varepsilon = \qty{2.53e-13}{s^{-1}}$.

\textbf{15.} $\eta = \num{62e6}/(2\times\num{1.27e-13}) = \qty{2.4e20}{Pa.s}$,
dicht bij de \qty{3e20}{Pa.s} van de bovenmantel onder het Lake Ardo: de zone
was even zwak als de mantel.

\textbf{16.} $e^{-E/RT}$ verandert met de factor 0.064 bij \qty{400}{\celsius}
(zestien keer zo langzaam) en met 11.0 bij \qty{500}{\celsius}.

\textbf{17.} $(78/62)^4 = 2.5$ keer zo snel.

\textbf{18.} $t = 9.51/\num{2.53e-13} = \qty{3.75e13}{s}$.

\textbf{19.} $\num{3.75e13}/\num{3.156e10} = \qty{1190}{kyr}$.

\textbf{20.} $\qty{3804}{m}/\qty{1.19e6}{yr} = \qty{3.2}{mm/yr}$.

\textbf{21.} Tien of meer keer zo langzaam als de belangrijkste plaatgrenzen: de zone
was een kleine tak, of een van meerdere die de beweging van een grens deelden.

\textbf{22.} $\dot\gamma = 0.02/(\num{3.156e7}\times400) = \qty{1.58e-12}{s^{-1}}$,
6.2 keer zo snel; de spanning stijgt met $6.2^{1/4} = 1.58$, tot \qty{98}{MPa}, en
de korrels zouden $3631/98^{1.26} = \qty{11}{\micro m}$ zijn.

\textbf{23.} De tijd verloopt als $\sigma^{-4}$: het korrelgroottebereik alleen geeft
$(62/78)^4 = 0.40$ tot $(62/52)^4 = 2.0$ keer de schatting, en de
onzekerheid van de exponent $1.26 \pm 0.13$ en van de constanten van de vloeiwet
verbreedt dat verder: het resultaat is een orde van grootte, geen datum.

\textbf{24.} De \omterm{def:b3:deformation-mechanisms:shear-zone}{Varden-schuifzone} kostte ongeveer
\textbf{\qty[text-series-to-math=true]{1190}{kyr}} nodig om zijn rek op te bouwen.
\end{solution}
